Inteligentné hodinky sa za posledné desaťročie stali bežnou
súčasťou každodenného života. Okrem zobrazovania notifikácií
a~sledovania zdravotných údajov ponúkajú aj prístup
k~pokročilým senzorom -- akcelerometru a~gyroskopu -- ktoré
umožňujú zachytávať pohyb zápästia v~reálnom čase. Tieto
možnosti sa čoraz častejšie využívajú v~športovom prostredí,
kde sa hľadajú spôsoby, ako uľahčiť prácu trénerov, hráčov aj
rozhodcov.

V~zápasení zohrávajú gestá rozhodcu kľúčovú úlohu pri
komunikácii rozhodnutí. Rozhodca na žinenke nimi signalizuje
udelené body, napomenutia, pasivitu či ukončenie zápasu
lopatkovým víťazstvom. Tieto signály sú nevyhnutné pre
transparentnosť priebehu zápasu, no zároveň prinášajú
priestor pre nepresnosti pri ich ručnom zaznamenávaní do
papierovej bodovacej karty. Práve kombinácia ustálených
pohybov ruky a~potreby presného zápisu robí zo zápasenia
vhodnú doménu pre automatizovanú detekciu gest pomocou
inteligentných hodiniek.

Cieľom tejto diplomovej práce je navrhnúť a~implementovať
systém na rozpoznávanie gest rozhodcu zápasenia pomocou
inteligentných hodiniek Amazfit T-Rex 3 s~operačným systémom
Zepp OS a~sprievodnej mobilnej aplikácie pre Android. Hodinky
nosené na ľavom zápästí snímajú dáta z~akcelerometra
a~gyroskopu a~odosielajú ich do Android aplikácie, kde
prebieha samotné rozpoznávanie. Rozpoznané gestá -- bodovanie,
napomenutie, pasivita a~TOUCHE -- sú následne zapisované do
priebehu zápasu, čím sa nahrádza ručný zápis. Navrhnuté gestá
vychádzajú z~oficiálnych pravidiel \Gls{uww} a~zo signálov,
ktoré rozhodcovia bežne používajú na žinenke.

V~prvej kapitole sa venujeme základným pojmom potrebným na
porozumenie problematike -- inteligentným hodinkám,
operačnému systému Zepp OS a~senzorom, ktoré sú v~hodinkách
zabudované. Zvláštnu pozornosť venujeme akcelerometru
a~gyroskopu, ktoré tvoria základ rozpoznávania gest.

Druhá kapitola sa zameriava na zápasenie ako šport. Popisuje
jeho hlavné štýly, priebeh zápasu, úlohu rozhodcu a~situácie,
v~ktorých rozhodca udeľuje body, napomenutia či pasivitu.
Osobitne sa venujeme významu gest rozhodcov a~ich úlohe pri
komunikácii rozhodnutí, keďže práve tieto gestá tvoria základ
pre návrh nášho riešenia.

V~tretej kapitole predstavujeme vývojové prostriedky použité
pri tvorbe systému -- JavaScript a~Zeus \Gls{cli} pre vývoj
aplikácie pre Zepp OS, Kotlin a~Android Studio pre mobilnú
aplikáciu, ako aj Python s~knižnicami Pandas a~Matplotlib
slúžiacimi na analýzu a~vizualizáciu nameraných dát zo
senzorov.

Štvrtá kapitola popisuje výber hardvéru -- inteligentných
hodiniek. Definujeme funkčné požiadavky na zariadenie
a~zdôvodňujeme voľbu modelu Amazfit T-Rex 3, ktorý spĺňa
všetky kritériá pre snímanie pohybových dát aj pre vývoj
vlastných aplikácií.

V~piatej kapitole analyzujeme existujúce riešenia z~oblasti
rozhodcovských aplikácií pre smart hodinky, aplikácií
zameraných na zápasenie a~akademických prác venovaných
rozpoznávaniu gest pomocou inerciálnych senzorov na zápästí.
Cieľom je zmapovať silné a~slabé stránky doterajších prístupov
a~identifikovať priestor pre vlastný príspevok.

Šiesta kapitola sa venuje návrhu gest. Definujeme elementárne
pohyby ruky -- Rise Arm, Hand Back, Flick a~Hand Down -- z~ktorých
sa skladajú kompozitné gestá zodpovedajúce signálom rozhodcu.
Popisujeme aj prepojenie jednotlivých gest s~príslušnými
udalosťami v~zápase.

V~siedmej kapitole popisujeme samotnú implementáciu systému --
aplikáciu pre Zepp OS, ktorá zbiera dáta zo senzorov, a~Android
aplikáciu, ktorá tieto dáta spracováva, vyhodnocuje gestá
pomocou nakalibrovaných prahových hodnôt a~zobrazuje stav
zápasu obsluhe. Súčasťou kapitoly je aj popis vibračnej
spätnej väzby a~vývojových režimov aplikácie.

Posledná, ôsma kapitola sa zaoberá testovaním. Overujeme
spoľahlivosť rozpoznávania jednotlivých gest pri ich
izolovanom vykonávaní aj pri kompozitných sekvenciách
zodpovedajúcich reálnym scenárom použitia. Výsledky
testovania slúžia na vyhodnotenie funkčnosti systému a~na
identifikáciu jeho silných a~slabých stránok.

V~závere zhrnieme dosiahnuté výsledky, prínos práce
a~načrtneme možnosti ďalšieho rozvoja systému. Veríme, že
navrhnuté riešenie môže prispieť k~zjednodušeniu práce
rozhodcov v~zápasení a~zároveň poslúžiť ako základ pre
podobné aplikácie v~iných športoch s~ustálenou gestikulačnou
komunikáciou.